\documentclass[10pt]{beamer}
\usepackage[utf8]{inputenc}
\usepackage[T1]{fontenc}
\usepackage[italian]{babel}
\usepackage{biblatex}
\usepackage{xcolor}
\usetheme{Berlin}
\usecolortheme{beaver}
\usepackage{amsfonts}
\usepackage{amsmath}
\usepackage{dsfont}
\usepackage{graphicx}
\usepackage{graphics}
\usepackage{tikz}
\usepackage{array}
\usepackage{multirow}
\usepackage{multicol}
\usepackage{changepage}
\usepackage{booktabs}
\setbeamertemplate{footline}[frame number]
\setbeamertemplate{navigation symbols}{}
\setbeamerfont{itemize/enumerate subbody}{size=\normalsize} 

\addbibresource{biblio.bib}

\title{\textbf{LSH Near Duplicate Document Detection}}
\date{}

\begin{document}
\begin{frame}[plain]
    \titlepage
    \vfill
    \begin{center}
        \vspace{0.5cm}
        Nicola Castelletti, Pietro Stangherlin
    \end{center}
\end{frame}



\section{Indice}
\begin{frame}{Indice}
\begin{itemize}
  \item Background
 \item Obiettivi
 \item Metodi utilizzati
 \item Esperimenti
 \item Conclusioni
\end{itemize}
\end{frame}

\section{Background}
\begin{frame}{Background}
\begin{itemize}
 \item Problema:in grandi raccolte di documenti (in particolare testuali) è possibile che vi siano dei documenti duplicati o quasi uguali; è inoltre possibile che tali documenti abbiano identificativi diversi o comunque non direttamente riconducibili tra loro.
 \item Esempi
 \begin{itemize}
 \item medesimi documenti scannerizzati più volte da entità diverse
 \item un’organizzazione che mantiene versioni diverse, ovvero con lievi modifiche, di uno stesso documento
 \end{itemize}
\end{itemize}
\end{frame}

\section{Obiettivi}
\begin{frame}{Obiettivi}
Implementare un metodo basato su MinHash e LSH per identificare i documenti duplicati e quasi-duplicati.\\

Valutarne l’efficacia e ed effettuare considerazioni sull’efficienza.
\end{frame}

\section{Metodi utilizzati}
\begin{frame}{Metodi utilizzati}
\begin{itemize}
	\item Creazione della collezione sperimentale
	\item Shingling
	\item MinHash
	\item LSH
\end{itemize}
\end{frame}

\section{Creazione della collezione sperimentale}
\begin{frame}{Creazione della collezione sperimentale}
\begin{enumerate}
	\item Si assume di avere una collezione iniziale senza duplicati e quasi – duplicati. Ogni documento presenta un codice identificativo (id).
	\item Si seleziona (ad esempio tramite un generatore di numeri casuali) un sottoinsieme di documenti da cui generare i duplicati e i quasi duplicati.
	\item Quasi duplicati: sostituzione di alcuni caratteri con altri (casuali), per simulare un risultato di un programma di riconoscimento ottico di testo. Il parametro relativo alla frequenza di sostituzione è la percentuale sul numero totale di caratteri nel testo
	\item A ogni quasi–duplicato si assegna un id univoco (id2) e l’id del documento da cui è stato generato (id3). Viene creata una nuova collezione con i quasi-duplicati e i documenti originali (per questi ultimi il campo l’id3 è posto uguale a None)
\end{enumerate}

\begin{center}
\resizebox{0.9\textwidth}{!}{
\begin{tikzpicture}
% First Table (left)
\node (table1) at (0,0) {
\begin{tabular}{|c|c|}
\hline
\textbf{Id} & \textbf{text} \\
\hline
1 & hello world \\
\hline
2 & Bob, Alice and Ulm \\
\hline
3 & nosce te ipsum \\
\hline
\end{tabular}
};

% Second Table (right)
\node (table2) at (8,0) {
\begin{tabular}{|c|c|c|}
\hline
\textbf{id2} & \textbf{id3} & \textbf{text} \\
\hline
1 & None & hello world \\
\hline
2 & None & Bob, Alice and Ulm \\
\hline
3 & None & nosce te ipsum \\
\hline
4 & 3 & n@sCe t3 ipsuN \\
\hline
\end{tabular}
};

% Arrow
\draw[->, thick] (2.5,0) -- (5,0);
\end{tikzpicture}
}
\end{center}
\end{frame}

\section{MinHash}
\begin{frame}{MinHash}
Dimensione delle shingles: w = 9 \citep{leskovec2014mining} (parametro fissato, seguendo le indicazioni sul libro di testo)\\
Salvataggio signatures
\begin{itemize}
	\item scrittura su disco (database sql) (scelta impiegata)
	\item memoria centrale (es. btree) (non considerata, possibile per piccole collezioni)
\end{itemize}
Stima della memoria occupata dalle signatures (ms):
$$\text{ms} =  \text{dim_int} \times \text{n_perm} \times \text{n_doc}$$
\end{frame}


\section{Conclusioni}
\begin{frame}{Conclusioni}
\end{frame}

\printbibliography


\end{document}